% -----------------------------------------------------------
\section{Book of Abraham}
% -----------------------------------------------------------
	\label{BOOKOFABRAHAM}
	
% % ----------------------
% \subsection{See also}
% % ----------------------

% % ----------------------
% \subsection{1828 American Dictionary of the English Language} % *1828
% % ----------------------
	% \label{BOOKOFABRAHAM-1828}

	% \begin{compactitem}
		% \item
	% \end{compactitem}

% % ----------------------
% \subsection{Oxford English Dictionary} % *OED
% % ----------------------
	% \label{BOOKOFABRAHAM-OED}

	% \begin{compactitem}
		% \item[]
	% \end{compactitem}

% % ----------------------
% \subsection{Scriptural Usage} % *SCRIPTURES
% % ----------------------
	% \label{BOOKOFABRAHAM-SCRIPTURES}

	% % -----
	% \subsubsection{Old Testament} % *OT
	% % -----
		% \label{BOOKOFABRAHAM-OT}
	
		% % --
		% \paragraph{KJV}
		% % --
			% \label{BOOKOFABRAHAM-OT-KJV}
	
			% \begin{tabular}{ll}
				% Total occurrences in the OT & 
			% \end{tabular}
			
			% \begin{compactdesc}
				% \item[]
			% \end{compactdesc}
			
		% % --
		% \paragraph{Hebrew Bible}
		% % --
			% \label{BOOKOFABRAHAM-HEBREW}
		
			% %
			% \subparagraph{\texorpdfstring{\he{sample word}}{}}
			% %
				% \label{PAGA} % whatever is in step bible without periods
			
				% \begin{compactdesc}
					% \item[HALOT , PDF ] \textbf{}
						% \begin{compactitem}
							% \item[1]
						% \end{compactitem}
					% \item[BDB , PDF ] \textbf{}
						% \begin{compactitem}
							% \item[1]
						% \end{compactitem}
					% \item[DCH PDF ] \textbf{}
						% \begin{compactitem}
							% \item[1]
						% \end{compactitem}
				% \end{compactdesc}
				
				% \begin{compactdesc}
					% \item[Some OT uses] \textbf{}
						% \begin{compactdesc}
							% \item[]
						% \end{compactdesc}
				% \end{compactdesc}
			
		% % --
		% \paragraph{Septuagint}
		% % --
			% \label{BOOKOFABRAHAM-LXX}
		
			% %
			% \subparagraph{\texorpdfstring{\gr{sample word}}{}}
			% %
		
	% % -----
	% \subsubsection{New Testament} % *NT
	% % -----
		% \label{BOOKOFABRAHAM-NT}
	
		% % --
		% \paragraph{KJV}
		% % --
			% \label{BOOKOFABRAHAM-NT-KJV}
			
	
			% \begin{tabular}{ll}
				% Total occurrences in the NT & 
			% \end{tabular}
			
			% \begin{compactdesc}
				% \item[]
			% \end{compactdesc}
			
		% % --
		% \paragraph{Greek New Testament} % *GREEK
		% % --
			% \label{BOOKOFABRAHAM-GREEK}
		
			% %
			% \subparagraph{\texorpdfstring{\gr{sample word}}{}}
			% %
				% \label{KATALLAGE}
			
				% \begin{compactdesc}
					% \item[BDAG , PDF ] \textbf{}
						% \begin{compactitem}
							% \item 
						% \end{compactitem}
					% \item[CGL , PDF ] \textbf{}
						% \begin{compactitem}
							% \item 
						% \end{compactitem}
					% \item[LSJ , PDF ] \textbf{}
						% \begin{compactitem}
							% \item 
						% \end{compactitem}
				% \end{compactdesc}
				
				% \begin{compactdesc}
					% \item[Some NT uses] \textbf{}
						% \begin{compactdesc}
							% \item[]
						% \end{compactdesc}
				% \end{compactdesc}
		
	% % -----
	% \subsubsection{Book of Mormon} % *BOM
	% % -----
		% \label{BOOKOFABRAHAM-BOM}
	
		% \begin{tabular}{ll}
			% Total occurrences in the BOM & 
		% \end{tabular}
		
		% \begin{compactdesc}
			% \item[]
		% \end{compactdesc}
		
	% % -----
	% \subsubsection{Doctrine and Covenants} % *DC
	% % -----
		% \label{BOOKOFABRAHAM-DC}
	
		% \begin{tabular}{ll}
			% Total occurrences in the DC & 
		% \end{tabular}
		
		% \begin{compactdesc}
			% \item[]
		% \end{compactdesc}
		
	% % -----
	% \subsubsection{Pearl of Great Price} % *PGP
	% % -----
		% \label{BOOKOFABRAHAM-PGP}
	
		% \begin{tabular}{ll}
			% Total occurrences in the PGP & 
		% \end{tabular}
		
		% \begin{compactdesc}
			% \item[]
		% \end{compactdesc}
		
% % ----------------------
% \subsection{Key Phrases} % *KEYPHRASES *PHRASES
% % ----------------------
	% \label{BOOKOFABRAHAM-PHRASES}

	% \begin{compactdesc}
		% \item[] \textbf{\#times} | list
			% % \begin{compactdesc}
				% % \item[Bible anchor verses] \phantom{.}
					% % \begin{compactdesc}
						% % \item[Romans 5:5] \gr{}
					% % \end{compactdesc}
				% % \item[Commentary] ---
			% % \end{compactdesc}
	% \end{compactdesc}
	
% % ----------------------
% \subsection{Bible Dictionaries} % *DICTIONARIES
% % ----------------------
	% \label{BOOKOFABRAHAM-BD}

% % ----------------------
% \subsection{Restoration Usage} % *RESTORATION
% % ----------------------
	% \label{BOOKOFABRAHAM-RESTORATION}

	% % -----
	% \subsubsection{Joseph Smith} % *JS
	% % -----
		% \label{BOOKOFABRAHAM-JS}
	
		% \begin{compactdesc}
			% \item[{\hyperref[KEY]{Date/Source}}]
		% \end{compactdesc}
	
	% % -----
	% \subsubsection{Brigham Young} % *BY
	% % -----
		% \label{BOOKOFABRAHAM-BY}
	
		% \begin{compactdesc}
			% \item[{\hyperref[KEY]{Date/Source}}]
		% \end{compactdesc}
	
% ----------------------
\subsection{Commentary and Studies} % *COMMENTARY *STUDIES
% ----------------------
	\label{BOOKOFABRAHAM-COMMENTARY}

	% -----
	\subsubsection{Key Points}
	% -----
		\label{BOOKOFABRAHAM-KEYS}
	
		\begin{compactitem}
			\item -Michael Chandler showed up in Kirtland with the mummies and scrolls on or right before July 3, 1835 (JSEP 1).
		-JS and the saints bought the materials in July 1835 (JSP-J1:57).
		-JS showed Chandler some preliminary translations before he left Kirtland, according to JSEP 1.
		-JS's journal entry for October 1, 1835, records that he and Oliver Cowdery and W. W. Phelps were working on the Egyptian alphabet (JSP-J1:67).
		-JS thought that the papyri contained at least three different stories: one of Abraham, one of Joseph of Egypt, and one of an Egyptian princess named Katumin (JSEP 2).
		-Champollion's discoveries weren't known to the British public until 1837 (NMKMH 170, note 1).
		-Champollion's correct deciphering of Egyptian was not reported in America until December 28, 1842, in the \textit{New York Herald} (JSEP 2, note 3).
			For more on Champollion, see also JSEP 14---his books on grammar and vocabulary were not published until 1836 and 1841, respectively, in Paris. There was an 1830 publication in Boston based on some of Champollion's work.
		-The church canonized the Book of Abraham in 1880 (JSEP 3).
		-Attacks on the Book of Abraham came as early as 1861, but the major assault began in 1912, when an episcopal bishop named F. S. Spalding published a book called \textit{Joseph Smith, Jr., as a Translator}. He commissioned professional Egyptologists to take a look at the translation and they all panned it (JSEP 4--5). Remember they're basing their opinion on the woodcuts, the explanations, anachronisms, and so on, since at this point they thought they papyri were no longer extant.
		-The papyri were recovered on November 27, 1967, when the Metropolitan Museum of Art gifted them back to the church (JSEP 6). They knew they were the same as the ones JS had because the illustrations on the papyri matched the woodcut facsimiles.
		-\textit{Dialogue} published a translation from some Egyptologists in 1968, and they said that the fragments were from the Book of Breathings and the Book of the Dead.
		-JSEP 11--68 is an essay by Marquardt on the history of the JS Egyptian papers, with a nice timeline as well.
		-The Book of Abraham was translated from 1:1--2:18 in 1835; at that point in the translation JS leaves off and doesn't return to it again until around March 1842, right before the full version is going to appear in the \textit{Times and Seasons}.
			On March 9 1842 JS writes a letter to Edward Hunter in which he says ``I am now very busily engaged in Translating, and therefore cannot give as much time to Public matters as I could wish'' (PWJS 550).
			On March 8 1842 JS's journal entry says, ``Commenced Translating from the Book of Abraham, for the 10 No of the Times and seasons'' (JSP-J2:42).
		-So you have to remember that 1:1--2:18 are pre-Kirtland temple dedication stuff, and the rest is just a couple months (or less) before the restoration of the endowment on May 4 1842.
			For example, JS's journal entry for 1 March 1842, about a week or a week and a half before he would have translated/dictated Abraham 3, says that he was preparing the woodcut of one of the facsimiles, and that later in the evening he was with ``the Twelve \& their wives at Elder [Wilford] Woodruff's. --where he [JS] explained many important principles in relation to progressive improvement. in the scale of inteligent existince'' (JSP-J1:39).
				For this date WW's journal says, ``My Birth Day I am 35 years of age this day as the spring enters, And thinking I could not Celebrate my Birth day more effectually than to meet my best friends I killed a Turkey made a feast \& invited some of my friends to attend. Some came \& some could not. The company consisted of Joseph the Seer \& Lady Elders H. C. Kimball \& Lady J. Taylor \& Lady \& Mother, Dr. W. Richards \& John E. Page. After supper the evening was mostly spent in hearing Joseph the Seer convers about the blessings of the kingdom of God much to our edification'' (WWJ 2:156).
		
	Some connections from my reading:
		-In the grammar \& alphabet the characters are understood in terms of five different degrees. In a letter by Oliver Cowdery to William Frye, 22 December 1835, Cowdery writes, speaking of the original discovery of the mummies: ``There were several hundred mummies in the same catacomb. About one hundred embalmed after the first order and deposited and placed in niches and two or three hundred after the second and third order, and laid upon the floor or bottom of the grand cavity, the two last orders of embalmed were so decayed that they could not be removed and only eleven out of the first, found in the niches.'' So five ``orders'' here too.
		\item JS is also studying Greek at the same time---see the J1 entry for 23 December 1835.
		\item JS is studying Hebrew at the very same time as well, in late 1835 and the school starts in 1836; see the journal entries around that time in J1. This is extremely important and helps us understand a great deal of what happens with the translation of the Book of Abraham.
			\begin{compactitem}
				\item See, for example, the beginning of the entry for 19 january 1836.
					Also, 26 jan 1836; a very important entry is 7 March 1836, where he discusses translating Genesis 17 from Hebrew (see also 8 march of the same year, where they do Genesis 22 and some of Exodus 3, in preparation for 1 and 2 Psalms); another very important one is 30 January 1836, where Seixas himself declares that the manuscripts are ``original beyond all doubt''; 
				\item For more on JS and languages, see Juvenile Instructor 27, 1892, May 15, page 302, where a man named John Hess recollects (many, many years later) that JS was studying Greek and Latin in the fall of 1838.
			\end{compactitem}
		\end{compactitem}
		
	% % -----
	% \subsubsection{Sample Study}
	% % -----
		% \label{BOOKOFABRAHAM-study}